% from main.tex
%%%%%%%%%%  scan idx 491  |  folio 478  %%%%%%%%%%
\end{tikzcd}
\]
L'homomorphisme $I/I^2\to J/J^2$ se déduit de l'homomorphisme : $IB'\to J$ et coïncide avec celui qui se déduit de $p:I\to J$; l'homomorphisme : $i^*(\Omega^1_{V/Y})\longrightarrow j^*(\Omega^1_{V'/Y})$ se déduit du carré commutatif (1). La propriété a) est donc évidente.

L'homomorphisme : $K/K^2\to j^*(\Omega^1_{V'/V})$ considéré dans b) peut-être défini comme suit: soit $\overline x\in K/K^2$, soit $x$ un élément de $J$ remontant $\overline x$ (il en existe car $K=J/IB'$). On associe à $\overline x$ l'image de $x$ par l'homomorphisme composé : $J\xrightarrow{d}\Omega^1_{B'/A}\otimes_{B'}C\longrightarrow\Omega^1_{B'/B}\otimes_{B'}C$. Par ailleurs, soit $H$ l'idéal d'augmentation de : $B'\otimes_B B'\longrightarrow B'$ ; le carré commutatif (2) donne un homomorphisme : $H\to K$, d'où un homomorphisme : $H/H^2\otimes_{B'}C\longrightarrow K/K^2$. Or $H/H^2=\Omega^1_{B'/B}$ ; par suite, il suffit de vérifier que les homomorphismes construits sont inverses l'un de l'autre, ce qui se fait sans difficulté.

\textbf{Définition 2.3.} Soit $f:X\to Y$ un morphisme lissifiable. Le complexe $L^{X/Y}_\bullet$, défini à isomorphisme canonique près dans la catégorie dérivée $D(X)$, est appelé complexe cotangent relatif de $X$ sur $Y$ (ou de $f$).

\textbf{Proposition 2.4.} Soit $f:X\to Y$ un morphisme d'intersection complète (1.1) lissifiable. Alors $L^{X/Y}_\bullet$ est parfait. Plus précisément, pour toute
